\section{Related work}
\label{sec:related}

\subsection{Satellite-navigation interference in transport surveillance}

ADS-B horizontal position and velocity are derived from the aircraft's GNSS
receiver, so interference propagates directly into the surveillance picture
\citep{khan2025adsbsurvey,pik2024detecting,pik2025gps}. AIS has the same
weakness: position, SOG and COG are GNSS-derived and broadcast over an
unauthenticated VHF link \citep{itu2014ais}. Most research addresses
\emph{detection}. In aviation, studies correlate ADS-B position gaps and
integrity indicators with notices to airmen \citep{pik2024detecting}. At
sea, recent work uses AIS as a sensor of opportunity to detect and localise
interference over wide areas \citep{park2026aisbased,winkel2026seaspoof},
and recurrent classifiers flag spoofed AIS points \citep{raj2025vessel}.
These detectors decide \emph{when} a track can no longer be trusted. Our
work starts at that point and addresses what the operator sees
\emph{after} the decision.

\subsection{Trajectory prediction and reconstruction}

Classical trackers such as Kalman and interacting-multiple-model filters
\citep{kalman1960new,ren2024improved} work well over short horizons. Under
interference, however, corrupted measurements are discarded rather than
smoothed, so the filter runs open loop and its error grows with the gap.
Autoregressive models behave similarly. Recurrent networks
\citep{pang2019lstm,hashemi2024bidirectional} capture non-linear motion but
depend on training coverage and drift without feedback. Recent reviews of
aircraft \citep{kwak2025evolution} and ship trajectory prediction
\citep{li2023ship,guo2026systematic} describe a mature but domain-siloed
literature dominated by recurrent, convolutional, attention-based and hybrid
networks. Among maritime models, AIS-ACNet \citep{shin2024aisacnet} couples
parallel dilated-convolution encoders for position and for auxiliary SOG/COG
targets and was evaluated on Port of Houston AIS. Because it predicts
forward only and has public code, it is the most suitable literature
baseline for our protocol. TrAISformer \citep{nguyen2024transformer}
discretises AIS states and predicts them with a transformer. Generic
forecasters such as DLinear \citep{zeng2023transformers}, PatchTST
\citep{nie2023timeseries}, iTransformer \citep{liu2024itransformer}, Informer
\citep{zhou2021informer}, Autoformer \citep{wu2021autoformer}, FEDformer
\citep{zhou2022fedformer} and TimesNet \citep{wu2023timesnet} provide strong
domain-independent baselines, but each must be trained on target data.

A separate line of work imputes ordinary AIS coverage gaps. MH-GIN
\citep{liu2026mhgin} uses multi-scale heterogeneous graphs. HABIT
\citep{spiliopoulos2026habit} uses historical hexagon transitions and path
search. VISTA \citep{liu2026vista} lets an LLM choose an imputation pattern
from a knowledge graph built from historical counts, which gives each repair
a readable provenance. These methods are the closest to ours in spirit, but
they interpolate between two known endpoints and address benign gaps. Our
setting is forward-only: the end of the blackout is unknown, and the
positions inside it are deliberately discarded.

\subsection{Large language models and retrieval for temporal data}

LLMs can forecast numeric series zero-shot when values are serialised as
text \citep{gruver2023large}, and they can be reprogrammed for forecasting
with a small trained adapter \citep{jin2024timellm}. Retrieval-augmented
generation grounds a parametric model in non-parametric memory
\citep{lewis2020retrieval}, and TS-RAG \citep{ning2025tsrag} shows that
retrieval improves zero-shot time-series foundation models. For flights,
LLMs have been fine-tuned on serialised ADS-B states \citep{luo2025large}
and assessed qualitatively for reconstruction \citep{zhang2024exploratory}.
None of these approaches encodes the constraints that make trajectory
reconstruction safe to deploy: operational phase, a trusted signal that is
physically independent of the corrupted one, and a hard bound on feasible
displacement. None of them has been shown to move between transport domains
without retraining.

\subsection{Positioning}

Table~\ref{tab:positioning} summarises where this work sits. Its
distinguishing combination is a post-detection, forward-only formulation, no
target-domain training, an explicit physical guard and a demonstrated
transfer across two physically distinct transport domains.

\begin{table}[t]
\centering
\caption{Positioning of this work relative to representative
trajectory-prediction and reconstruction approaches. ``Post-detection''
means that horizontal positions after the onset are withheld rather than
denoised. ``No training'' means that no model parameters are fitted on
target-domain trajectories.}
\label{tab:positioning}
\footnotesize
\begin{tabular}{p{3.0cm}ccccc}
\toprule
Approach & Post-det. & No training & Phys. guard & Multi-domain & Interpretable \\
\midrule
Kalman / IMM filters & \checkmark & \checkmark & -- & \checkmark & -- \\
Recurrent/conv. predictors & -- & -- & -- & -- & -- \\
Generic time-series forecasters & -- & -- & -- & -- & -- \\
AIS imputation (MH-GIN, VISTA, HABIT) & -- & partly & -- & -- & partly \\
Fine-tuned LLM flight prediction & -- & -- & -- & -- & \checkmark \\
Zero-shot/retrieval LLM forecasting & -- & \checkmark & -- & \checkmark & -- \\
TrajRAG, aviation only \citep{felendler2026trajrag} & \checkmark & \checkmark & \checkmark & -- & \checkmark \\
\textbf{This work} & \checkmark & \checkmark & \checkmark & \checkmark & \checkmark \\
\bottomrule
\end{tabular}
\end{table}
